目标函数、调整费用线性化关系、母线平衡、储电量递推及边界约束均为线性关系，因此每个更新时刻的剩余时域问题均可写为稀疏线性规划。使用 HiGHS 求解时，先最小化式\eqref{eq:q3-objective}；在目标值的货币容差内，再最小化充放电吞吐量，以消除数值上可能出现的等价循环。该次级目标不会改变经济目标值，但能使调度轨迹更易解释。算法不包含非线性迭代或复杂分支，故以步骤描述其信息流和锁定关系即可，无须另绘流程图。

\noindent\textbf{四时点滚动购电调整算法}
\algorithmstart
\algostep{1}{构造日初信息集}
在每日 $0{:}00$ 汇集全日电价、日初储电量、此前已结束日期的负荷历史与当期光伏预报；按式\eqref{eq:q3-load-forecast}生成日前负荷预测，并将逐小时光伏预报按式\eqref{eq:q3-pv-interpolation}映射至 10 分钟时段。滚动回测中，预测训练样本严格早于决策时点\cite{tashman}。
\algostep{2}{生成日初承诺}
以预测供需、储能约束和终端价值求解日初 LP，得到 $g^0_{d,t}$ 及初始储能预案；该承诺成为当天所有调整结算的唯一基准。
\algostep{3}{实时执行一个时段}
在每个 10 分钟时段，仅读取当前实际负荷、实际光伏和上期储电量；在承诺上限 $x_{d,t}\le g_{d,t}$ 下求得充放电、弃光和必要的紧急购电，并更新储电量。
\algostep{4}{更新预测并锁定历史}
到达 $6{:}00$、$12{:}00$ 或 $18{:}00$ 时，以新的已发布光伏预报重建剩余时段预测；将已执行时段的购电承诺、物理状态和费用设为常数，并从式\eqref{eq:q3-first-mutable}规定的下一时段开始优化。
\algostep{5}{计算信息价值}
在锁定原剩余承诺和允许修改剩余承诺两种条件下分别求解 LP，计算 $\operatorname{VoI}_{\tau}$。若其大于 $0.01$ 元，则用自由解替换未执行部分；否则保留原承诺。
\algostep{6}{继续滚动直至日末}
重复 Step 3--Step 5，得到最终有效购电承诺、实际储能轨迹和紧急购电记录；跨日运行时，将日末储电量传递为下一日的日初状态。
\algostep{7}{核验并按题设输出}
核验能量平衡、储电量递推、功率边界、锁定时段和同时充放电；将日初计划、最终调整计划、四小时充放电汇总和紧急购电量按题目模板写入相应工作表。
\algorithmend

正式运行以 2025 年 1 月 1 日 $6000\,\mathrm{kWh}$ 为初始 SOC，全年连续传递状态；主口径采用因果负荷预测、线性锚点光伏映射、相对日初计划的一次性调整结算和 $48$ 小时终端价值。题设工作簿只导出 2--12 月，但策略优劣以全年连续 SOC 的实际成本比较，不以局部试算代替年度结论。
